% HAND-WRITTEN fragment -- not generated by maestro2tex, but placed by it via the
% `order` list in configs/tb_anatop_pixel_dacR2R12_corners.json.
%
% Numbers here come from tab_dacr2r12_mc_ideal, tab_dacr2r12_mc_scale and
% tab_dacr2r12_mc_wpsu, and from the corner table tab_dacr2r12_wpsu. Re-read them
% from those tables if either run is repeated.
%
% Source, ideal-rail figures and tables (rewritten 2026-09-06): the F16
% bit-driver campaign of 2026-09-05, 200 samples, seed 12345, LDS sampling,
% variations=all, section castalia_pm_25, 25 C, with the custom 2.5 V driver
% castalia/anatop_and2_dac -- simarchive/tb_anatop_pixel_dacR2R12/
% 20260905_driverfix/dac_mc/. Verified before use: 82 devices at mismatchflag=1
% and none at 0, 234 random variables, section castalia_pm_25 on the run that
% executed. Two earlier Monte Carlo attempts on this block failed one or other
% of those checks and produced tidy distributions carrying no matching
% information; their numbers must not be reintroduced.
%
% Source, tab_dacr2r12_mc_wpsu and everything else referenced from here:
% Plan.2.Run.3, the earlier campaign, on the 1.0 V standard-cell driver and
% through the supply block deleted from anatop_pixel on 2026-08-24.

\subsubsection{Matching}\label{sss:dacr2r12-mc}

R-2R linearity is a matching property, so it is the Monte Carlo run and not the corner
set that measures it. On an ideal \SI{2.5}{\volt} rail, over 200 samples varying process
and mismatch together with the custom \SI{2.5}{\volt} bit driver
\texttt{anatop\_and2\_dac} the tapeout cell carries, INL is \SI{1.49}{LSB} mean
against a \SI{1}{LSB} limit and the worst sample reaches \SI{3.39}{LSB}; DNL reaches
\SI{-6.50}{LSB} (Table~\ref{tab:dacr2r12-mc-ideal}). \SI{20.5}{\percent} of samples
meet \(\mathrm{INL} < \SI{1}{LSB}\) and \SI{42.0}{\percent} are monotonic
(\(\mathrm{DNL} > \SI{-1}{LSB}\)); \SI{18.5}{\percent} meet both, which over four
independent channels is a \SI{0.12}{\percent} chip yield.

The failure is matching and nothing else. Full scale, LSB and gain error are set by the
rail and by ladder attenuation, not by device ratios, and they barely move across the
same 200 samples: \(\sigma\) is \SI{1.2}{\micro\volt} on a \SI{2.4994}{\volt} full
scale and \SI{0.50}{ppm} about a \SI{-0.51}{ppm} gain error
(Table~\ref{tab:dacr2r12-mc-scale}).
All of the spread in Table~\ref{tab:dacr2r12-mc-ideal} is therefore resistor mismatch in
the ladder.

Referencing the ladder to \texttt{VddDac} instead of an ideal rail did not change this
either. The ratiometric figures of Table~\ref{tab:dacr2r12-mc-wpsu}
(\SI{1.51}{LSB} mean INL, \SI{-1.75}{LSB} mean DNL, \SI{40.0}{\percent} monotonic)
reproduce the ideal-rail figures to within \SI{5}{\percent}, because dividing the
reference out removes the supply block and leaves the same ladder. That table and the raw
columns beside it, \SI{36.8}{LSB} and \SI{-10.8}{LSB}, are the earlier campaign: the
1.0\,V bit driver, and the supply block of Section~\ref{sss:dacr2r12-supply} that the
tapeout channel no longer has. They are kept as the corroboration of the matching result,
not as a specification.

The corner set sees none of it. Ratiometric INL spans \SIrange{0.432}{0.450}{LSB} across
ff, fs, sf and ss, a \SI{4}{\percent} spread, while the Monte Carlo samples of the
same quantity scatter from \SI{0.35}{LSB} to \SI{3.49}{LSB}. A process corner moves every
ladder resistor in the same direction, which a ratio cancels; the corner table is context
and not a linearity result.

INL and DNL are both skewed (Fisher skew \num{+0.74} and \num{-1.14}), so the tables
quote the median, the first and ninety-ninth percentiles and the worst sample rather than
mean \(\pm\,3\sigma\).
